%==============================================================================
\section{Conclusion and Future Work}
\label{sec:conclusion}

We placed an information-theoretic measurement layer on an existing accelerated-
aging burn-in platform, re-expressing environmental contamination, PVT
correction and aging recoverability as entropy, mutual information and channel
capacity. On a \valDurationHours~h PID campaign we established the empirical
anchors ($\Ent{\slk}=\valHs$~bits; linear $R^2=\valRsqLinear$) and a methodology
that replaces the model-dependent $R^2$ with mutual information, defines burn-in
quality as the minimization of $\MI{\slk}{\Tdie,\Vcore}$, converts measured noise
into resolvable aging states and an alarm-channel capacity, and bounds RUL
detectability honestly via Cram\'er--Rao and Bayesian estimation. A
Wiener-process/Kalman-filter extension of the same estimation toolbox further
turns this single-point RUL bound into a continuously updated, denoised aging
trajectory and a closed-form Inverse-Gaussian RUL \emph{distribution}
(\cref{sec:res-wiener}), cross-validated against the Cram\'er--Rao/Bayesian
estimate and against the MDL model-selection result of \cref{sec:res-correction}.

\paragraph{Future work.} Acquire a matched-length bang-bang campaign to complete
the cross-regime KL and capacity comparison on fresh data; extend to multiple
devices to separate a permanent-aging trajectory from device-to-device
variation; deploy the windowed pipeline online as a continuous information-feature
generator feeding the SLM controller, with the Kalman-smoothed trajectory of
\cref{sec:res-wiener} as its live aging estimate; calibrate an actual failure
threshold $D$ (e.g.\ from a timing-margin specification) so
\cref{eq:invgauss} yields an operational RUL distribution rather than the
self-referential thresholds used here; and pursue the ICA/blind-separation
route with additional derived channels to strengthen model-free aging
recovery.

\section*{Reproducibility}
Every scalar in this paper is generated by the analysis pipeline into
\texttt{generated/values.tex} and every figure/table into \texttt{figures/} and
\texttt{tables/}; re-running the analysis on a new log regenerates the document
without manual transcription.

\section*{Acknowledgments}
This work was developed for the Information Theory course (TIP7266, PPGETI/UFC,
2026.1) and builds on the author's research on integrated-circuit aging sensors
and burn-in instrumentation.
